\slajdid{09_2-s055}
\begin{frame}[label=09_2-s055]
\gusttekst\setlength{\abovedisplayskip}{3pt}\setlength{\belowdisplayskip}{3pt}\setlength{\abovedisplayshortskip}{2pt}\setlength{\belowdisplayshortskip}{2pt}\begin{columns}[T,onlytextwidth]\column{.37\textwidth}\centering\begin{tikzpicture}[x=.55cm,y=.6cm,font=\sitneoznake,>=Latex]
\ctikzset{bipoles/length=.6cm}
\draw(0,1.2)--(0,1.8);\draw(-.13,1.15)--(-.13,1.85);\draw(-.23,1.5)circle(.1);
\draw(-.8,1.5)node[ocirc]{}node[below,align=center]{$V_I$\\$=0$}--(-.33,1.5);
\draw(0,1.75)--(.35,1.75)--(.35,2.9);\draw(.1,2.9)--(.6,2.9)node[above left]{$V_{DD}$};
\draw(0,1.25)--(.35,1.25)--(.35,.9)--(1.4,.9)node[ocirc]{};
\node[below,align=center]at(.6,-1.1){$V_O$\\$(0\to V_{DD})$};
\draw(1.1,.9)to[C,l=$C_L$](1.1,-.65)node[ground]{};
\begin{scope}[shift={(3.8,0)}]
\draw(0,1.2)--(0,1.8);\draw(-.13,1.15)--(-.13,1.85);
\draw(-.8,1.5)node[ocirc]{}node[below,align=center]{$V_I$\\$=V_{DD}$}--(-.13,1.5);
\draw(0,1.75)--(.35,1.75)--(.35,2.25)--(1.3,2.25)node[ocirc]{};
\node[below,align=center]at(.6,-1.1){$V_O$\\$(V_{DD}\to0)$};
\draw(0,1.25)--(.35,1.25)--(.35,-.65)node[ground]{};
\draw(1.15,2.25)to[C,l=$C_L$](1.15,-.65)node[ground]{};
\end{scope}
\end{tikzpicture}
\par\bigskip\raggedright
U toku pražnjenja kondenzatora celokupna energija koja se nalazi u njemu će biti disipirana na N tranzistoru.
\[E_N=E_{CL}=\frac{C_LV_{DD}^2}2\]
\column{.61\textwidth}
U situaciji kada P tranzistor treba da napuni parazitnu kapacitivnost ukupna energija koja će se potrošiti iz izvora za napajanje je
\[\begin{aligned}E_{VDD}&=\int_0^\infty i_{VDD}(t)V_{DD}\,dt=V_{DD}\int_0^\infty i_{CL}(t)dt\\&=V_{DD}\int_0^\infty C_L\frac{u_{CL}(t)}{dt}dt\end{aligned}\]
\[E_{VDD}=C_LV_{DD}\int_0^{V_{DD}}du_{CL}(t)=C_LV_{DD}^2\]
dok je na kraju procesa punjenja energija u kondenzatoru
\[E_{CL}=\frac{C_LV_{DD}^2}2\]
Razlika energija je disipirana na P tranzistoru.
\[E_P=E_{VDD}-E_{CL}=C_LV_{DD}^2-\frac{C_LV_{DD}^2}2=\frac{C_LV_{DD}^2}2\]
\end{columns}
\end{frame}
